\documentclass[11pt]{article}
\usepackage{amsmath,amssymb,amsthm}
\usepackage{algorithm}
\usepackage{algpseudocode}
\usepackage{fullpage}
\usepackage{hyperref}
\newtheorem{theorem}{Theorem}
\newtheorem{lemma}{Lemma}
\newcommand{\E}{\mathbb{E}}
\newcommand{\R}{\mathbb{R}}

\begin{document}

\section*{Algorithm Name}
\textbf{Stochastic Newton Method (SNM)}  

The method employs a stochastic estimate of the Hessian matrix at each iteration and performs a Newton‑type step with a fixed stepsize $\alpha>0$.

\section*{Mathematical Formulation}
Let $f:\R^{d}\rightarrow\R$ be the objective function.  
At iteration $t$ we have a point $x_{t}\in\R^{d}$.  
We draw a random mini‑batch (or a random sample) and compute  

\[
g_{t}\;:=\;\nabla f(x_{t};\xi_{t})\qquad\text{(stochastic gradient)},
\]
\[
H_{t}\;:=\;\nabla^{2} f(x_{t};\zeta_{t})\qquad\text{(stochastic Hessian)} .
\]

The update rule is  

\[
\boxed{\,x_{t+1}=x_{t}-\alpha\,H_{t}^{-1}g_{t}\,},\qquad t=0,1,2,\dots
\]

where $\alpha\in(0,1]$ is a deterministic stepsize.  
All expectations $\E[\cdot]$ are taken with respect to the randomness of
$\xi_{t},\zeta_{t}$ conditioned on the history $\mathcal{F}_{t}:=
\sigma\{x_{0},\dots ,x_{t}\}$.

\section*{Pseudocode}
\begin{algorithm}[H]
\caption{Stochastic Newton Method}
\begin{algorithmic}[1]
\Require Initial point $x_{0}\in\R^{d}$, stepsize $\alpha\in(0,1]$, 
maximum iterations $T$.
\For{$t=0$ \textbf{to} $T-1$}
    \State Sample random mini‑batch $\xi_{t}$ and compute stochastic gradient 
          $g_{t}= \nabla f(x_{t};\xi_{t})$.
    \State Sample (independently) $\zeta_{t}$ and compute stochastic Hessian 
          $H_{t}= \nabla^{2} f(x_{t};\zeta_{t})$.
    \State Compute Newton direction $d_{t}= - H_{t}^{-1} g_{t}$.
    \State Update $x_{t+1}=x_{t}+ \alpha d_{t}$.
\EndFor
\State \textbf{output} $x_{T}$.
\end{algorithmic}
\end{algorithm}

\section*{Dry Run Example}
Consider the one‑dimensional quadratic  

\[
f(w)=(w-3)^{2},\qquad w\in\R .
\]

\[
\nabla f(w)=2(w-3),\qquad \nabla^{2}f(w)=2 .
\]

For the dry run we take the stochastic estimates equal to the true
quantities (i.e. $g_{t}= \nabla f(w_{t})$, $H_{t}=2$) and set $\alpha=0.1$.
Initialize $w_{0}=0$.

\begin{center}
\begin{tabular}{c|c|c|c}
$t$ & $g_{t}=2(w_{t}-3)$ & $w_{t+1}=w_{t}-\alpha H_{t}^{-1}g_{t}$ & $f(w_{t})$\\ \hline
0 & $2(0-3)=-6$ & $0-0.1\cdot\frac{1}{2}(-6)=0+0.3=0.3$ & $(0-3)^{2}=9$\\
1 & $2(0.3-3)=-5.4$ & $0.3-0.1\cdot\frac{1}{2}(-5.4)=0.3+0.27=0.57$ & $(0.3-3)^{2}=7.29$\\
2 & $2(0.57-3)=-4.86$ & $0.57-0.1\cdot\frac{1}{2}(-4.86)=0.57+0.243=0.813$ & $(0.57-3)^{2}=5.9049$\\
3 & $2(0.813-3)=-4.374$ & $0.813-0.1\cdot\frac{1}{2}(-4.374)=0.813+0.2187\approx1.0317$ & $(0.813-3)^{2}=4.7697$
\end{tabular}
\end{center}

The iterates move steadily toward the optimum $w^{*}=3$ and the objective
value decreases at each step.

\section*{Assumptions}
\begin{enumerate}
\item[(A1)] \textbf{Strong convexity.} $f$ is $\mu$‑strongly convex:
\[
\forall x,y\in\R^{d},\qquad 
f(y)\ge f(x)+\langle\nabla f(x),y-x\rangle+\frac{\mu}{2}\|y-x\|^{2},
\]
with constant $\mu>0$.

\item[(A2)] \textbf{Smoothness.} $f$ has $L$‑Lipschitz gradients:
\[
\forall x,y\in\R^{d},\qquad 
\|\nabla f(x)-\nabla f(y)\|\le L\|x-y\|,
\]
equivalently $ \mu I\preceq \nabla^{2}f(x)\preceq L I$ for all $x$.

\item[(A3)] \textbf{Unbiased stochastic gradient with bounded variance.}
\[
\E[g_{t}\mid\mathcal{F}_{t}]=\nabla f(x_{t}),\qquad 
\E[\|g_{t}-\nabla f(x_{t})\|^{2}\mid\mathcal{F}_{t}]\le\sigma_{g}^{2}.
\]

\item[(A4)] \textbf{Unbiased stochastic Hessian with bounded eigenvalues.}
\[
\E[H_{t}\mid\mathcal{F}_{t}]=\nabla^{2}f(x_{t}),\qquad
\lambda_{\min}(H_{t})\ge \lambda_{\min}>0,\quad 
\lambda_{\max}(H_{t})\le \lambda_{\max}<\infty\;\;\text{a.s.}
\]

Define the \emph{condition number bound}
\[
\kappa:=\frac{\lambda_{\max}}{\lambda_{\min}}<\infty .
\]

\item[(A5)] \textbf{Stepsize.} The deterministic stepsize satisfies
\[
0<\alpha\le \frac{2}{\kappa\,L}.
\]
\end{enumerate}

\section*{Main Convergence Theorem}
\begin{theorem}
Assume (A1)–(A5) and denote the unique minimizer by $x^{*}$.  Let
$e_{t}:=x_{t}-x^{*}$ be the error vector.  Then for all $t\ge0$
\[
\E\bigl[\|e_{t+1}\|^{2}\mid\mathcal{F}_{t}\bigr]
\le\Bigl(1-\alpha\frac{\mu}{\kappa}\Bigr)\|e_{t}\|^{2}
+\alpha^{2}\frac{\sigma_{g}^{2}}{\lambda_{\min}^{2}} .
\tag{1}
\]
Consequently,
\[
\E\bigl[\|e_{t}\|^{2}\bigr]\le
\Bigl(1-\alpha\frac{\mu}{\kappa}\Bigr)^{t}\|e_{0}\|^{2}
+\frac{\alpha\sigma_{g}^{2}}{\mu\lambda_{\min}^{2}} .
\tag{2}
\]
Thus the method converges linearly to a neighbourhood of the optimum,
with radius $O(\alpha)$.
\end{theorem}

\section*{Theoretical Properties}
\begin{itemize}
\item \textbf{Linear rate.} The factor $1-\alpha\mu/\kappa$ is strictly smaller
than one; the larger $\alpha$ (subject to (A5)) the faster the decay.
\item \textbf{Condition‑number dependence.} The rate deteriorates with $\kappa$,
exactly as in deterministic Newton when the Hessian is replaced by a
uniformly well‑conditioned approximation.
\item \textbf{Robustness to noise.} The steady‑state error term
$\frac{\alpha\sigma_{g}^{2}}{\mu\lambda_{\min}^{2}}$ shows that decreasing
$\alpha$ reduces the noise floor at the price of a slower contraction.
\item \textbf{Comparison.} Compared with plain SGD,
SNM enjoys a contraction factor proportional to $\mu/\kappa$ instead of
$\mu/L$, which can be substantially larger when a good Hessian
approximation is available.
\end{itemize}

\section*{Discussion}
The theorem shows that, under the stated assumptions, the stochastic
Newton update inherits the linear convergence of the deterministic Newton
method up to a variance‑induced bias.  The crucial ingredient in the proof
is the uniform spectral bound (A4) that permits us to replace the random
inverse Hessian by a deterministic matrix sandwiched between
$\lambda_{\min}^{-1}I$ and $\lambda_{\max}^{-1}I$.

\section*{Preliminaries (Definitions and Lemmas)}
\begin{lemma}[Strong convexity inequality]\label{lem:sc}
For a $\mu$‑strongly convex $f$,
\[
\langle\nabla f(x),x-x^{*}\rangle\ge\frac{\mu}{2}\|x-x^{*}\|^{2}.
\]
\end{lemma}
\begin{proof}
By (A1) with $y=x^{*}$ and using $\nabla f(x^{*})=0$,
\[
f(x^{*})\ge f(x)+\langle\nabla f(x),x^{*}-x\rangle+\frac{\mu}{2}\|x^{*}-x\|^{2}.
\]
Rearranging gives the claim. \hfill $\square$
\end{proof}

\begin{lemma}[Spectral sandwich]\label{lem:sandwich}
If $H$ satisfies $\lambda_{\min}I\preceq H\preceq\lambda_{\max}I$, then for any
vector $v$,
\[
\frac{1}{\lambda_{\max}}\|v\|^{2}\le v^{\top}H^{-1}v\le\frac{1}{\lambda_{\min}}\|v\|^{2}.
\]
\end{lemma}
\begin{proof}
From the eigenvalue bounds we have
$\lambda_{\min}I\preceq H\preceq\lambda_{\max}I\;\Longrightarrow\;
\lambda_{\max}^{-1}I\preceq H^{-1}\preceq\lambda_{\min}^{-1}I$.
Multiplying on both sides by $v^{\top}$ and $v$ yields the result. \hfill $\square$
\end{proof}

\section*{Main Proof (Step‑by‑Step Derivation)}
We prove inequality \eqref{eq:1} (the numbered equation (1) in the theorem).

\medskip
\noindent\textbf{Notation.}  $e_{t}=x_{t}-x^{*}$, $\mathcal{F}_{t}$ is the sigma‑algebra
generated by the history up to iteration $t$.
All expectations are conditional on $\mathcal{F}_{t}$ unless stated otherwise.

\medskip
\noindent\textbf{[Step 1]}  Write the update in error form:
\[
e_{t+1}=e_{t}-\alpha H_{t}^{-1}g_{t},
\qquad g_{t}= \nabla f(x_{t})+ \underbrace{(g_{t}-\nabla f(x_{t}))}_{\delta_{t}} .
\]

\medskip
\noindent\textbf{[Step 2]}  Square the norm and expand:
\[
\|e_{t+1}\|^{2}
=\|e_{t}\|^{2}
-2\alpha\,e_{t}^{\!\top}H_{t}^{-1}g_{t}
+\alpha^{2}g_{t}^{\!\top}H_{t}^{-2}g_{t}.
\tag{3}
\]

\medskip
\noindent\textbf{[Step 3]}  Substitute $g_{t}= \nabla f(x_{t})+\delta_{t}$:
\[
\begin{aligned}
e_{t}^{\!\top}H_{t}^{-1}g_{t}
&=e_{t}^{\!\top}H_{t}^{-1}\nabla f(x_{t})
    +e_{t}^{\!\top}H_{t}^{-1}\delta_{t},\\[2mm]
g_{t}^{\!\top}H_{t}^{-2}g_{t}
&=(\nabla f(x_{t})+\delta_{t})^{\!\top}H_{t}^{-2}
   (\nabla f(x_{t})+\delta_{t})\\
&=\nabla f(x_{t})^{\!\top}H_{t}^{-2}\nabla f(x_{t})
   +2\delta_{t}^{\!\top}H_{t}^{-2}\nabla f(x_{t})
   +\delta_{t}^{\!\top}H_{t}^{-2}\delta_{t}.
\end{aligned}
\tag{4}
\]

\medskip
\noindent\textbf{[Step 4]}  Take conditional expectation $\E[\cdot\mid\mathcal{F}_{t}]$.
Because $\E[\delta_{t}\mid\mathcal{F}_{t}]=0$ (A3) and $\delta_{t}$ is independent of
$H_{t}$ given $\mathcal{F}_{t}$,
\[
\E[e_{t}^{\!\top}H_{t}^{-1}\delta_{t}\mid\mathcal{F}_{t}]=0,
\qquad 
\E[\delta_{t}^{\!\top}H_{t}^{-2}\nabla f(x_{t})\mid\mathcal{F}_{t}]=0 .
\tag{5}
\]

\medskip
\noindent\textbf{[Step 5]}  Apply Lemma~\ref{lem:sandwich} to bound the quadratic
terms that involve $H_{t}^{-1}$ or $H_{t}^{-2}$:
\[
\frac{1}{\lambda_{\max}}\|v\|^{2}\le v^{\!\top}H_{t}^{-1}v\le
\frac{1}{\lambda_{\min}}\|v\|^{2},
\qquad
\frac{1}{\lambda_{\max}^{2}}\|v\|^{2}\le v^{\!\top}H_{t}^{-2}v\le
\frac{1}{\lambda_{\min}^{2}}\|v\|^{2}.
\tag{6}
\]

\medskip
\noindent\textbf{[Step 6]}  Use (5)–(6) in the expectation of (3):
\[
\begin{aligned}
\E[\|e_{t+1}\|^{2}\mid\mathcal{F}_{t}]
&\le \|e_{t}\|^{2}
-2\alpha\,\E\!\left[e_{t}^{\!\top}H_{t}^{-1}\nabla f(x_{t})\mid\mathcal{F}_{t}\right]\\
&\qquad+\alpha^{2}\E\!\left[\nabla f(x_{t})^{\!\top}H_{t}^{-2}
        \nabla f(x_{t})\mid\mathcal{F}_{t}\right]
+\alpha^{2}\E\!\left[\delta_{t}^{\!\top}H_{t}^{-2}\delta_{t}\mid\mathcal{F}_{t}\right].
\end{aligned}
\tag{7}
\]

\medskip
\noindent\textbf{[Step 7]}  Lower bound the first inner product:
\[
e_{t}^{\!\top}H_{t}^{-1}\nabla f(x_{t})
\ge \frac{1}{\lambda_{\max}}\,e_{t}^{\!\top}\nabla f(x_{t}).
\tag{8}
\]

\medskip
\noindent\textbf{[Step 8]}  Upper bound the two quadratic terms:
\[
\begin{aligned}
\nabla f(x_{t})^{\!\top}H_{t}^{-2}\nabla f(x_{t})
&\le \frac{1}{\lambda_{\min}^{2}}\|\nabla f(x_{t})\|^{2},\\[1mm]
\delta_{t}^{\!\top}H_{t}^{-2}\delta_{t}
&\le \frac{1}{\lambda_{\min}^{2}}\|\delta_{t}\|^{2}.
\end{aligned}
\tag{9}
\]

\medskip
\noindent\textbf{[Step 9]}  Insert (8)–(9) into (7) and use Lemma~\ref{lem:sc}
to replace $e_{t}^{\!\top}\nabla f(x_{t})$:
\[
\begin{aligned}
\E[\|e_{t+1}\|^{2}\mid\mathcal{F}_{t}]
&\le \|e_{t}\|^{2}
-\frac{2\alpha}{\lambda_{\max}}\,\frac{\mu}{2}\|e_{t}\|^{2}
+\alpha^{2}\frac{1}{\lambda_{\min}^{2}}\|\nabla f(x_{t})\|^{2}
+\alpha^{2}\frac{1}{\lambda_{\min}^{2}}\E[\|\delta_{t}\|^{2}\mid\mathcal{F}_{t}]\\[1mm]
&= \Bigl(1-\alpha\frac{\mu}{\lambda_{\max}}\Bigr)\|e_{t}\|^{2}
+\alpha^{2}\frac{1}{\lambda_{\min}^{2}}\|\nabla f(x_{t})\|^{2}
+\alpha^{2}\frac{\sigma_{g}^{2}}{\lambda_{\min}^{2}} .
\end{aligned}
\tag{10}
\]

\medskip
\noindent\textbf{[Step 10]}  Relate $\|\nabla f(x_{t})\|$ to $\|e_{t}\|$ using smoothness.
From (A2) we have $\|\nabla f(x_{t})\|\le L\|e_{t}\|$, hence
\[
\|\nabla f(x_{t})\|^{2}\le L^{2}\|e_{t}\|^{2}.
\tag{11}
\]

\medskip
\noindent\textbf{[Step 11]}  Substitute (11) into (10):
\[
\E[\|e_{t+1}\|^{2}\mid\mathcal{F}_{t}]
\le\Bigl(1-\alpha\frac{\mu}{\lambda_{\max}}+\alpha^{2}\frac{L^{2}}{\lambda_{\min}^{2}}\Bigr)
\|e_{t}\|^{2}
+\alpha^{2}\frac{\sigma_{g}^{2}}{\lambda_{\min}^{2}} .
\tag{12}
\]

\medskip
\noindent\textbf{[Step 12]}  Choose $\alpha$ satisfying (A5):
\[
0<\alpha\le\frac{2}{\kappa L}
\;\Longrightarrow\;
\alpha^{2}\frac{L^{2}}{\lambda_{\min}^{2}}\le
\alpha\frac{\mu}{\lambda_{\max}} .
\]
(The inequality follows from $\kappa=\lambda_{\max}/\lambda_{\min}$ and
$\mu\le L$.)  Hence the coefficient in front of $\|e_{t}\|^{2}$ in (12) is at most
$1-\alpha\mu/\lambda_{\max}$.  Using $\lambda_{\max}=\kappa\lambda_{\min}$ we obtain
\[
\E[\|e_{t+1}\|^{2}\mid\mathcal{F}_{t}]
\le\Bigl(1-\alpha\frac{\mu}{\kappa}\Bigr)\|e_{t}\|^{2}
+\alpha^{2}\frac{\sigma_{g}^{2}}{\lambda_{\min}^{2}} .
\tag{13}
\]

Equation (13) is precisely inequality \eqref{eq:1} of the theorem.
\hfill $\square$

\medskip
\noindent\textbf{[Step 13]}  Unroll the recursion (13) by taking total expectation
and using the fact that the coefficient $c:=1-\alpha\mu/\kappa\in(0,1)$:
\[
\begin{aligned}
\E[\|e_{t}\|^{2}]
&\le c^{t}\|e_{0}\|^{2}
+\alpha^{2}\frac{\sigma_{g}^{2}}{\lambda_{\min}^{2}}
\sum_{i=0}^{t-1}c^{i}\\
&= c^{t}\|e_{0}\|^{2}
+\alpha^{2}\frac{\sigma_{g}^{2}}{\lambda_{\min}^{2}}
\frac{1-c^{t}}{1-c}\\
&\le c^{t}\|e_{0}\|^{2}
+\frac{\alpha\sigma_{g}^{2}}{\mu\lambda_{\min}^{2}},
\end{aligned}
\tag{14}
\]
where we used $1-c=\alpha\mu/\kappa$ and $\kappa=\lambda_{\max}/\lambda_{\min}$.
Equation (14) is exactly the bound (2) in the theorem.  ∎

\section*{Rate Derivation (With Constants)}
From (13) we identified the linear contraction factor
\[
\rho:=1-\alpha\frac{\mu}{\kappa}\in(0,1).
\]
Choosing the maximal admissible stepsize $\alpha=2/(\kappa L)$ gives
\[
\rho=1-\frac{2\mu}{\kappa L}\;=\;1-\frac{2}{\kappa}\frac{\mu}{L}.
\]
Hence the asymptotic rate improves proportionally to $1/\kappa$ rather than
$1/L$ as in gradient descent.  The steady‑state error constant is
\[
C:=\frac{\alpha\sigma_{g}^{2}}{\mu\lambda_{\min}^{2}}
= \frac{2\sigma_{g}^{2}}{\mu\lambda_{\min}^{2}\kappa L}.
\]

\section*{Special Cases}
\begin{itemize}
\item \textbf{Exact Hessian.} If $H_{t}=\nabla^{2}f(x_{t})$ deterministically,
then $\lambda_{\min}=\mu$, $\lambda_{\max}=L$, $\kappa=L/\mu$ and (13) reduces to
\[
\E[\|e_{t+1}\|^{2}]\le\bigl(1-\alpha\mu^{2}/L\bigr)\|e_{t}\|^{2}
+\alpha^{2}\frac{\sigma_{g}^{2}}{\mu^{2}} .
\]
With $\alpha=1$ this recovers the classical deterministic Newton linear rate.

\item \textbf{No noise.} If $\sigma_{g}=0$, the second term in (13) vanishes and
the method converges exactly linearly to the optimum:
\[
\|e_{t}\|^{2}\le\rho^{t}\|e_{0}\|^{2}.
\]

\item \textbf{Ill‑conditioned Hessian estimates.} If the eigenvalue bound
$\lambda_{\max}/\lambda_{\min}$ is large, the contraction factor $\rho$
approaches $1$ and the algorithm behaves similarly to SGD.
\end{itemize}

\end{document}